%% LyX 2.4.1 created this file.  For more info, see https://www.lyx.org/.
%% Do not edit unless you really know what you are doing.
\documentclass[english]{article}
\usepackage[T1]{fontenc}
\usepackage[latin9]{inputenc}
\usepackage{geometry}
\geometry{verbose,tmargin=1in,bmargin=1in,lmargin=1in,rmargin=1in,headheight=1in,headsep=1in,footskip=1in}
\usepackage{amsmath}

\makeatletter
%%%%%%%%%%%%%%%%%%%%%%%%%%%%%% User specified LaTeX commands.
\date{Due at the beginning of class on Thu Jun 24th}
% Added by lyx2lyx
\setlength{\parskip}{\smallskipamount}
\setlength{\parindent}{0pt}

\makeatother

\usepackage{babel}
\begin{document}
\title{Problem Set 1}
\maketitle
\begin{center}
    Answers must be typed. Send solutions to raffaele.saggio@ubc.ca. You can work in teams under the constraint $\max(Team Size)=2$
\par\end{center}


\subsection*{Question 1: Returns to Schooling}

Consider an infinitely-lived person making a choice about how long
to go to school. She faces a discount rate $r$. If she chooses $S$
years of schooling, she earns $y(S)$ from date $S$ onward. While
at school, she earns nothing in the labor market, but receives a scholarship
payment. Specifically, if she attends $S$ years of schooling, she
is paid a transfer $T$ from time $0$ to time $S$. Her goal is to
maximize the present value of earnings and scholarship payments.
\begin{enumerate}
\item Write down a continuous-time expression for the person's objective
function. Derive and interpret a condition to characterize the optimal
duration of schooling.
\item Suppose that agents differ in ability, $\alpha_{i}$, and the human
capital production function for individual $i$ takes the form
\begin{equation}
y_{i}(S)=\exp(\alpha_{i}+\beta S)
\end{equation}

Solve for the optimal duration of schooling. Do higher-ability agents
get more or less schooling? Why?

\item Now suppose both tuition payments and ability
differ across agents: the scholarship offered agent $i$ is $T_{i}$.
Derive the probability limit of the slope coefficient from a bivariate
OLS regression of log earnings on schooling in terms of $\beta$,
$\sigma_{T}^{2}=Var(\log T_{i})$, $\sigma_{\alpha}^{2}=Var(\alpha_{i})$,
and $\sigma_{T\alpha}=Cov(\log T_{i},\alpha_{i})$. Discuss how the
bias of OLS varies with these parameters.

\item Now suppose that $T_{i}=T$ (tuition payments are constant in the population) but $y_{i}(S)=\exp(\alpha_{i}+\beta_{i} S)$ (proportional returns to schooling differ across individuals). Derive the probability limit of the slope coefficient from a bivariate
OLS regression of log earnings on schooling in terms of $\beta$,
$\sigma_{\beta}^{2}=Var(\beta_{i})$, $\sigma_{\alpha}^{2}=Var(\alpha_{i})$,
and $\sigma_{\beta,\alpha}=Cov(\beta_{i},\alpha_{i})$. 


\end{enumerate}


\subsection*{Question 2: Keane and Wolpin (1997)}

This question asks you to solve and estimate a simplified version
of the model from Keane and Wolpin (1997). An agent makes choices
at two ages, $a\in\{1,2\}$. At each age, she chooses between three
alternatives: working in a white collar occupation, working in a blue
collar occupation, and going to school. These alternatives are labeled
$m\in\{1,2,3\}$. For alternatives 1 and 2, the flow utility of working
at age $a$ is

\begin{center}
$R_{m}(a)=\alpha_{m}+\beta_{m}G(a)+\gamma_{m}X_{m}(a)+\epsilon_{m}(a)$,
\par\end{center}

where $G(a)$ is education and $X_{m}(a)$ is experience in occupation
$m$. The utility of going to school is

\begin{center}
$R_{3}(a)=\epsilon_{3}(a)$.
\par\end{center}

The agent enters period 1 with no experience or education. Assume
that the $\epsilon_{m}(a)$ are iid draws from an extreme value type
I distribution, independent over time. The $\epsilon_{m}(2)$ are
unknown to the agent in period 1. The agent has a discount rate $\delta\in(0,1)$.
\begin{enumerate}
\item Let $S(2)$ denote the vector of state variables entering age $2$
(before the realization of $\epsilon_{m}(2)$). What are the elements
of this vector? What is the probability that the agent chooses each
alternative in period 2, as a function of $S(2)$?
\item Define the value function in period 2, 
\end{enumerate}
\begin{center}
$V_{2}(S(2))={\displaystyle \max_{m\in\{1,2,3\}}R_{m}(2;S(2))}.$
\par\end{center}

Write down an expression for $E\left[V_{2}(S(2))\right]$ for each
possible value of $S(2)$ prior to the realization of $\epsilon_{m}(2)$.
(Hint: In a discrete choice problem where the utility of alternative
$k$ is $U_{ik}+\eta_{ik}$, where the $\eta_{ik}$ are distributed
iid extreme value type I, we have $E\left[\max_{k}U_{ik}+\eta_{ik}\right]=\nu+\log\left(\sum_{k}\exp(U_{ik})\right)$,
where $\nu$ is a constant of integration that you can ignore.)
\begin{enumerate}
\item [3.]\setcounter{enumi}{3} Use your results from part (2) to write
down an expression for the agent's expected payoff from each choice
in period 1. What is the \emph{ex ante} probability of choosing each
alternative in period 1, prior to the realization of $\epsilon_{m}(1)$?
\item [4.]\setcounter{enumi}{4} Suppose you have data for a sample of agents
choosing according to this model. Let $C_{i}(a)\in\{1,2,3\}$ denote
individual $i$'s choice at age $a$. The observed data for individual
$i$ is the choice sequence $(C_{i}(1),C_{i}(2))$. Use your results
from parts (1) and (3) to write down the likelihood that individual
$i$ chooses the sequence $(C_{i}(1),C_{i}(2))$. Take the log and
sum over individuals to obtain an expression for the sample log-likelihood.
\item [5.]\setcounter{enumi}{5} Download the data set ``kw\_example.csv''
from the course website. This data set includes 5,000 simulated observations
from the model. Import the data into a statistical software package
of your choosing. Estimate the parameter vector
\end{enumerate}
\begin{center}
$\theta\equiv(\alpha_{1},\beta_{1},\gamma_{1},\alpha_{2},\beta_{2},\gamma_{2},\delta)^{\prime}$
\par\end{center}

by maximum likelihood. (I recommend using Matlab and minimizing the
negative of the log-likelihood using the \emph{fminunc} routine. You
may wish to search over log $\delta$ rather than $\delta$ to keep
the discount rate positive at every evaluation of the likelihood.)
Attach your computer code, and make a table showing your parameter
estimates; you may also provide standard errors for extra credit.
Briefly provide some economic interpretation for your parameter estimates
and discuss the tradeoffs faced by the agent when choosing an alternative
in period 1.
\end{document}
